\chapter{Segmentation}

\section{Introduction}

Segmentation is a memory management technique that divides a program into logical units called \textbf{segments}. Unlike paging, which divides memory into fixed-size blocks, segmentation reflects the programmer's view of memory and groups related data together.

Typical program components naturally map to segments:

\begin{itemize}
\item Code Segment
\item Data Segment
\item Stack Segment
\item Heap Segment
\item Shared Library Segment
\end{itemize}

Segmentation improves program organization, protection, and sharing.

\section{Motivation for Segmentation}

\subsection{Problem with Linear Address Spaces}

Consider a program containing:

\begin{itemize}
\item Instructions
\item Global Variables
\item Functions
\item Stack
\item Dynamic Memory
\end{itemize}

Treating all of these as a single contiguous memory block is inefficient.

Different parts of a program require different:

\begin{itemize}
\item Protection Levels
\item Access Rights
\item Growth Patterns
\item Sharing Policies
\end{itemize}

\subsection{Logical Organization}

Programmers think of memory as:

\begin{center}
Code + Data + Stack + Heap
\end{center}

not as one large array of bytes.

Segmentation aligns memory management with the programmer's perspective.

\section{Segmentation Fundamentals}

\subsection{Definition}

Segmentation divides a process into variable-sized logical units called segments.

Each segment contains logically related information.

\subsection{Examples}

\begin{table}[h]
\centering
\begin{tabular}{|c|l|}
\hline
Segment Number & Contents \
\hline
0 & Code \
1 & Global Data \
2 & Heap \
3 & Stack \
4 & Shared Libraries \
\hline
\end{tabular}
\caption{Example Segments}
\end{table}

\subsection{Logical Address Format}

In segmentation, a logical address consists of:

[
\langle Segment\ Number,\ Offset \rangle
]

Example:

[
\langle 2, 150 \rangle
]

means:

\begin{itemize}
\item Segment Number = 2
\item Offset = 150
\end{itemize}

\section{User View vs Physical Memory View}

\subsection{User View}

Programmers see:

\begin{itemize}
\item Functions
\item Variables
\item Stack
\item Heap
\item Libraries
\end{itemize}

\subsection{Physical Memory View}

Hardware sees:

\begin{itemize}
\item Physical Memory Addresses
\item Memory Frames
\item Actual RAM Locations
\end{itemize}

\subsection{Illustration}

\begin{figure}[h]
\centering
\begin{tikzpicture}

\node[draw,minimum width=4cm,minimum height=5cm] (user) {};
\node at (0,2) {Code};
\node at (0,1) {Data};
\node at (0,0) {Heap};
\node at (0,-1) {Stack};

\node[draw,minimum width=4cm,minimum height=5cm,right=5cm of user] (phys) {};

\node at (5,2) {Physical Block A};
\node at (5,1) {Physical Block B};
\node at (5,0) {Physical Block C};
\node at (5,-1) {Physical Block D};

\draw[->] (user) -- (phys);

\end{tikzpicture}
\caption{User View vs Physical Memory View}
\end{figure}

\section{Segment Structure}

Each segment contains:

\begin{itemize}
\item Base Address
\item Length (Limit)
\item Protection Information
\end{itemize}

\subsection{Example}

\begin{table}[h]
\centering
\begin{tabular}{|c|c|c|}
\hline
Segment & Base & Limit \
\hline
0 & 1000 & 500 \
1 & 4000 & 800 \
2 & 9000 & 300 \
\hline
\end{tabular}
\caption{Segment Structure Example}
\end{table}

\section{Segment Tables}

\subsection{Definition}

A Segment Table stores information about all segments belonging to a process.

\subsection{Fields}

Typically contains:

\begin{itemize}
\item Segment Number
\item Base Address
\item Limit
\item Protection Bits
\end{itemize}

\subsection{Example Segment Table}

\begin{table}[h]
\centering
\begin{tabular}{|c|c|c|}
\hline
Segment & Base & Limit \
\hline
0 & 1000 & 600 \
1 & 4000 & 800 \
2 & 7000 & 500 \
3 & 9000 & 1000 \
\hline
\end{tabular}
\caption{Example Segment Table}
\end{table}

\section{Address Translation in Segmentation}

\subsection{Translation Process}

Logical Address:

[
\langle s,d \rangle
]

where:

\begin{itemize}
\item (s) = segment number
\item (d) = offset
\end{itemize}

Steps:

\begin{enumerate}
\item Locate segment table entry.
\item Check offset against limit.
\item Add offset to base address.
\item Generate physical address.
\end{enumerate}

\subsection{Formula}

If

[
d < Limit
]

then

[
Physical\ Address = Base + d
]

\subsection{Numerical Example}

Segment Table:

\begin{center}
Segment 1: Base = 4000, Limit = 800
\end{center}

Logical Address:

[
\langle 1,200 \rangle
]

Physical Address:

[
4000 + 200 = 4200
]

\section{Address Translation Diagram}

\begin{figure}[h]
\centering
\begin{tikzpicture}

\node[draw] (cpu) {CPU};

\node[draw,right=3cm of cpu] (st) {Segment Table};

\node[draw,right=3cm of st] (mem) {Physical Memory};

\draw[->] (cpu) -- node[above]{Segment, Offset} (st);

\draw[->] (st) -- node[above]{Base + Offset} (mem);

\end{tikzpicture}
\caption{Segmentation Address Translation}
\end{figure}

\section{Protection Mechanisms}

Segmentation naturally supports protection.

\subsection{Protection Bits}

Each segment may have:

\begin{itemize}
\item Read Permission
\item Write Permission
\item Execute Permission
\end{itemize}

\subsection{Example}

\begin{table}[h]
\centering
\begin{tabular}{|c|c|}
\hline
Segment & Permission \
\hline
Code & Execute Only \
Data & Read/Write \
Stack & Read/Write \
\hline
\end{tabular}
\caption{Segment Protection}
\end{table}

\section{Sharing Mechanisms}

Segmentation allows efficient sharing.

\subsection{Example}

Two processes may share:

\begin{itemize}
\item Code Segment
\item Shared Libraries
\end{itemize}

while maintaining separate:

\begin{itemize}
\item Data Segments
\item Stack Segments
\end{itemize}

\subsection{Benefits}

\begin{itemize}
\item Reduced Memory Usage
\item Improved Efficiency
\end{itemize}

\section{Segmentation Hardware}

Hardware support typically includes:

\begin{itemize}
\item Segment Table Base Register (STBR)
\item Segment Table Length Register (STLR)
\item MMU Support
\end{itemize}

\subsection{STBR}

Points to segment table location.

\subsection{STLR}

Stores number of segments.

\section{External Fragmentation}

\subsection{Problem}

Segments have variable sizes.

Memory develops holes over time.

\subsection{Example}

\begin{verbatim}
| P1 | Free | P2 | Free | P3 |
\end{verbatim}

Although sufficient free memory exists, large requests may fail.

\subsection{Result}

External fragmentation occurs.

\section{Compaction}

\subsection{Definition}

Compaction moves segments to create contiguous free space.

\subsection{Before}

\begin{verbatim}
| P1 | Free | P2 | Free | P3 |
\end{verbatim}

\subsection{After}

\begin{verbatim}
| P1 | P2 | P3 | Large Free Block |
\end{verbatim}

\subsection{Drawback}

Compaction is expensive.

\section{Pure Segmentation}

\subsection{Definition}

Memory management using only segmentation.

\subsection{Characteristics}

\begin{itemize}
\item Variable-sized segments
\item Logical organization
\item External fragmentation
\end{itemize}

\section{Segmentation with Paging}

\subsection{Motivation}

Segmentation suffers from external fragmentation.

Paging solves fragmentation problems.

\subsection{Approach}

Each segment is divided into pages.

\subsection{Benefits}

\begin{itemize}
\item Logical Organization
\item Reduced Fragmentation
\item Better Memory Utilization
\end{itemize}

\subsection{Architecture}

\begin{figure}[h]
\centering
\begin{tikzpicture}

\node[draw] (seg) {Segment};

\node[draw,right=3cm of seg] (page1) {Page 1};

\node[draw,right=5cm of seg] (page2) {Page 2};

\node[draw,right=7cm of seg] (page3) {Page 3};

\draw[->] (seg) -- (page1);
\draw[->] (seg) -- (page2);
\draw[->] (seg) -- (page3);

\end{tikzpicture}
\caption{Segmentation with Paging}
\end{figure}

\section{Intel x86 Segmentation Overview}

\subsection{Historical Design}

Intel x86 originally relied heavily on segmentation.

Logical address:

[
Segment:Offset
]

Example:

[
CS:IP
]

[
DS:Offset
]

\subsection{Segment Registers}

\begin{itemize}
\item CS (Code Segment)
\item DS (Data Segment)
\item SS (Stack Segment)
\item ES
\item FS
\item GS
\end{itemize}

\subsection{Modern x86}

Modern operating systems use:

\begin{itemize}
\item Flat Address Spaces
\item Paging Dominantly
\item Minimal Segmentation
\end{itemize}

\section{Modern Operating Systems and Segmentation}

\subsection{Linux}

Modern Linux primarily uses:

\begin{itemize}
\item Paging
\item Virtual Memory
\end{itemize}

Segmentation usage is minimal.

\subsection{Windows}

Modern Windows also uses:

\begin{itemize}
\item Paging-Based Memory Management
\item Flat Memory Model
\end{itemize}

\subsection{Reason}

Paging provides:

\begin{itemize}
\item Better Flexibility
\item Reduced Fragmentation
\item Easier Memory Management
\end{itemize}

\section{Advantages of Segmentation}

\begin{enumerate}
\item Matches Programmer View
\item Natural Protection Support
\item Easy Sharing
\item Logical Program Organization
\item Independent Segment Growth
\end{enumerate}

\section{Disadvantages of Segmentation}

\begin{enumerate}
\item External Fragmentation
\item Compaction Overhead
\item Complex Allocation
\item Variable Size Management
\item Additional Hardware Support
\end{enumerate}

\section{Real-World Use Cases}

\subsection{Program Organization}

Separate:

\begin{itemize}
\item Code
\item Data
\item Stack
\end{itemize}

\subsection{Shared Libraries}

Shared code segments among processes.

\subsection{Security}

Different permissions per segment.

\subsection{Legacy x86 Systems}

Used segmentation extensively.

\section{Segmentation vs Paging}

\begin{table}[h]
\centering
\begin{tabular}{|p{4cm}|p{5cm}|p{5cm}|}
\hline
Feature & Segmentation & Paging \
\hline
Division Basis & Logical Units & Fixed Size Pages \
\hline
Size & Variable & Fixed \
\hline
Programmer Visible & Yes & No \
\hline
Fragmentation & External & Internal \
\hline
Protection & Natural & Supported via Page Tables \
\hline
Sharing & Easy & Possible \
\hline
Allocation & Complex & Simpler \
\hline
Modern Usage & Limited & Dominant \
\hline
\end{tabular}
\caption{Segmentation vs Paging}
\end{table}

\section{Key Takeaways}

\begin{itemize}
\item Segmentation divides memory into logical units.
\item Logical addresses contain segment number and offset.
\item Segment tables store base and limit values.
\item Segmentation supports protection naturally.
\item Segmentation supports sharing efficiently.
\item External fragmentation is a major drawback.
\item Compaction reduces fragmentation.
\item Modern systems primarily use paging.
\item Segmentation with paging combines strengths of both.
\item Intel x86 historically used segmentation extensively.
\end{itemize}

\section{20 Interview Questions with Answers}

\begin{enumerate}
\item What is segmentation?\
Answer: A memory management technique dividing memory into logical segments.

\item Why was segmentation introduced?\
Answer: To match programmer view of memory.

\item What is a segment?\
Answer: A logical unit of a program.

\item What is a logical address in segmentation?\
Answer: Segment number and offset.

\item What is a segment table?\
Answer: A table storing segment information.

\item What does a segment table contain?\
Answer: Base address and limit.

\item What is STBR?\
Answer: Segment Table Base Register.

\item What is STLR?\
Answer: Segment Table Length Register.

\item How is physical address calculated?\
Answer: Base + Offset.

\item What happens if offset exceeds limit?\
Answer: Trap or protection fault.

\item Does segmentation support protection?\
Answer: Yes.

\item Does segmentation support sharing?\
Answer: Yes.

\item What fragmentation occurs in segmentation?\
Answer: External fragmentation.

\item What is compaction?\
Answer: Rearranging memory to create large free blocks.

\item What is pure segmentation?\
Answer: Segmentation without paging.

\item Why combine segmentation with paging?\
Answer: To reduce fragmentation.

\item What segment register stores code segment?\
Answer: CS.

\item What segment register stores stack segment?\
Answer: SS.

\item Do modern Linux systems rely heavily on segmentation?\
Answer: No.

\item Which memory technique dominates modern OSes?\
Answer: Paging.
\end{enumerate}

\section{20 MCQs with Answers}

\begin{enumerate}
\item Segmentation divides memory into:
A) Frames
B) Pages
C) Segments
D) Blocks

Answer: C

\item Segment sizes are:
A) Fixed
B) Variable
C) Equal
D) Random

Answer: B

\item Logical address contains:
A) Page + Frame
B) Segment + Offset
C) Base + Limit
D) Frame + Offset

Answer: B

\item Segment table stores:
A) Base and Limit
B) Frames
C) Pages
D) Cache Lines

Answer: A

\item STBR stands for:
Answer: Segment Table Base Register

\item STLR stands for:
Answer: Segment Table Length Register

\item Segmentation supports:
A) Protection
B) Sharing
C) Both
D) Neither

Answer: C

\item Physical address equals:
A) Base + Offset
B) Page + Offset
C) Limit + Offset
D) Segment + Offset

Answer: A

\item Main fragmentation issue:
A) Internal
B) External
C) Cache
D) TLB

Answer: B

\item Compaction addresses:
A) Internal Fragmentation
B) External Fragmentation
C) TLB Misses
D) Cache Misses

Answer: B

\item Segment register for code:
A) DS
B) SS
C) CS
D) ES

Answer: C

\item Segment register for stack:
A) CS
B) SS
C) DS
D) FS

Answer: B

\item Segmentation is:
A) Programmer Visible
B) Invisible
C) Kernel Only
D) Hardware Only

Answer: A

\item Paging uses:
A) Variable Sizes
B) Fixed Sizes
C) Logical Units
D) Functions

Answer: B

\item Segmentation uses:
A) Logical Units
B) Fixed Blocks
C) Cache Lines
D) Frames Only

Answer: A

\item Modern Linux primarily uses:
A) Segmentation
B) Paging
C) Swapping Only
D) Overlaying

Answer: B

\item Modern Windows primarily uses:
A) Segmentation
B) Paging
C) Compaction
D) Overlays

Answer: B

\item Shared libraries benefit from:
A) Segmentation
B) Cache
C) DMA
D) TLB

Answer: A

\item Segment table entry includes:
A) Limit
B) Base
C) Protection
D) All

Answer: D

\item Segmentation with paging combines:
A) Benefits of both techniques
B) Only paging
C) Only segmentation
D) Neither

Answer: A
\end{enumerate}

\section{Revision Notes}

\begin{itemize}
\item Segmentation divides memory into logical units.
\item Logical address = Segment Number + Offset.
\item Segment Table stores Base and Limit.
\item MMU performs address translation.
\item Protection and sharing are natural in segmentation.
\item External fragmentation is the primary problem.
\item Compaction reduces fragmentation.
\item Pure segmentation uses variable-sized segments.
\item Segmentation with paging is more efficient.
\item Modern operating systems primarily rely on paging.
\end{itemize}

\section{Cheat Sheet}

\begin{center}
\begin{tabular}{|p{5cm}|p{8cm}|}
\hline
Logical Address & Segment Number + Offset \
\hline
Physical Address & Base + Offset \
\hline
Segment Table & Base + Limit + Protection \
\hline
STBR & Segment Table Base Register \
\hline
STLR & Segment Table Length Register \
\hline
Major Advantage & Logical Organization \
\hline
Major Disadvantage & External Fragmentation \
\hline
Compaction & Removes External Fragmentation \
\hline
Protection & Natural Support \
\hline
Sharing & Easy Support \
\hline
Pure Segmentation & Variable Sized Segments \
\hline
Segmentation + Paging & Hybrid Approach \
\hline
Intel x86 & Historically Segmentation Based \
\hline
Modern Linux & Paging Dominant \
\hline
Modern Windows & Paging Dominant \
\hline
\end{tabular}
\end{center}
